\noindent Proposto por \cite{vapnik1997support}, o algoritmo
de Máquinas de Vetores de Suporte (SVM) permite construir
fronteiras de decisão não lineares por meio da indução de um
hiperplano linear em um espaço de características de alta dimensão,
obtido por meio de uma transformação do espaço original. Este capítulo
aborda o SVM nos contextos de classificação, regressão e ranking. A
apresentação do caso de classificação tem como objetivo servir de base
para os demais, uma vez que estes herdam seus principais conceitos. Por
fim, o capítulo se restringe à descrição dos casos não separáveis.


\section{Máquinas de Vetores de Suporte para Classificação}

O objetivo do SVM para classificação é encontrar um hiperplano
que maximize a margem de separação entre as classes. Formalmente,
conforme definido por \citeonline[p.~79]{mohri2018foundations},
considere um espaço de covariáveis $\mathcal{X} \subseteq \mathbb{R}^d$
e um espaço de rótulos $\mathcal{Y} = {-1, +1}$. Seja $\mathit{f} : \mathcal{X} \rightarrow \mathcal{Y}$
a função alvo desconhecida.

Dado um conjunto de hipóteses $\mathcal{H}$ de funções que
mapeiam $\mathcal{X}$ em $\mathcal{Y}$, o problema de classificação
binária consiste em encontrar uma hipótese $\mathit{h} \in \mathcal{H}$
que minimize o erro de generalização:
$$
\mathit{R}_{\mathcal{D}}\left(\mathit{h}\right) = \underset{\mathit{x} \sim \mathcal{D}}{\mathbb{P}}\left[\mathit{h}(\mathit{x}) \neq \mathit{f}(\mathit{x})\right]\text{,}
$$
em que $\mathcal{D}$ é uma distribuição desconhecida sobre $\mathcal{X}$.

Para isso, o modelo recebe um conjunto de treinamento
$\mathit{S} = \left((\mathit{x}_1, \mathit{y}_1), \ldots, (\mathit{x}_m, \mathit{y}_m)\right) \in (\mathcal{X} \times \mathcal{Y})^m$
composto por $\mathit{m}$ observações i.i.d. de $\mathcal{X} \times \mathcal{Y}$,
geradas segundo a distribuição $\mathcal{D}$, tais que $\mathit{y}_i = \mathit{f}(\mathit{x}_i)$.
Um conjunto de hipóteses natural é um hiperplano, definido da seguinte forma:

\begin{figure}[htb]
 \caption{Dois hiperplanos de separação. O hiperplano à direita é aquele que maximiza a margem de separação.}
 \label{fig:jabref}
 \centering
 \includegraphics[scale=0.45]{hyperplane.png}
Fonte: \citeonline[p.~80]{mohri2018foundations}
\end{figure}
$$
\mathcal{H} = \{\mathbf{x} \mapsto \sgn(\mathbf{w} \cdot \mathbf{x} + b) : \mathbf{w} \in \mathbb{R}^d, b \in \mathbb{R}\}
$$
em que $\mathbf{w} \cdot \mathbf{x} + b = 0$ é a equação de um hiperplano
em $\mathbb{R}^d$, com $\mathbf{w} \in \mathbb{R}^d$ um vetor normal não nulo
ao hiperplano e $b \in \mathbb{R}$ um escalar. Dessa forma, a hipótese
$\mathbf{x} \mapsto \sgn(\mathbf{w} \cdot \mathbf{x} + b)$ define a regra
de classificação
$$
\mathit{h}\left(\mathbf{x}\right) = \sgn\left(\mathbf{w} \cdot \mathbf{x} + b\right)
$$
como ilustrado na Figura \ref{fig:jabref}.

Observe que, na Figura \ref{fig:jabref}, podem existir infinitos hiperplanos
que separam as classes. No entanto, conforme discutido anteriormente, a solução do SVM
baseia-se na escolha do hiperplano que maximiza a margem geométrica entre as classes.

A margem geométrica $\rho_{\mathit{h}}\left(\mathbf{x}\right)$ de um classificador
$\mathit{h}: \mathbf{x} \mapsto \sgn(\mathbf{w} \cdot \mathbf{x} + b)$, em um ponto $\mathbf{x}$,
é definida como a sua distância Euclidiana ao hiperplano $\mathbf{w} \cdot \mathbf{x} + b = 0$:
$$
\rho_{\mathit{h}}\left(\mathit{x}\right) = \frac{|\mathbf{w} \cdot \mathbf{x} + b|}{\|\mathbf{w}\|_{2}}\text{.}
$$
Como a margem deve ser garantida para todas as observações, para uma amostra
$\mathit{S} = \left(\mathbf{x}_1, \ldots, \mathbf{x}_m\right)$, define-se a margem
geométrica do classificador $\mathit{h}$ como o menor valor entre as margens dos pontos
da amostra valor entre as margens dos pontos da amostra, isto é,
$$
\rho_{\mathit{h}} = \min_{i \in \left[\mathit{m}\right]} \rho_{\mathit{h}} \left(\mathbf{x}_i\right)\text{,}
$$
que corresponde aos pontos da amostra mais próximos ao hiperplano.
